\section{Timeline of Substantive Modifications}

This timeline selects changes that altered the architecture, precedence, temporal cycle, universal inventory, or United States overlay. It does not list routine reprints, every canonization, or every local grant.

\begin{longtable}{>{\raggedright\arraybackslash}p{0.13\linewidth}>{\raggedright\arraybackslash}p{0.25\linewidth}>{\raggedright\arraybackslash}p{0.54\linewidth}}
\toprule
Date & Act or development & Substantive effect\\
\midrule
\endhead
Early centuries & Sunday, Pascha, and martyrs' anniversaries & Established the weekly and annual Paschal centers and the local anniversary principle from which the temporal and sanctoral cycles grew. Precise beginnings are not reducible to a single decree.\\
Fourth--sixth centuries & Nativity--Epiphany and mature Paschal seasons & Consolidated a second temporal center and the preparatory and festal seasons surrounding Christmas and Easter, with continued regional variation.\\
1568--1570 & Roman Breviary and Missal after Trent & Codified and widely extended the Roman calendar and its books while preserving qualifying ancient local uses.\\
1911 & Pius X, \emph{Divino afflatu} & Reworked the Psalter and precedence so that festal accumulation would no longer displace the ferial cursus in the same way.\\
1951--1956 & Pius XII's Paschal Vigil, Holy Week, and simplification reforms & Restored the principal Holy Week rites to their proper hours, revised their structure, and suppressed many vigils and octaves while simplifying calendar ranks.\\
1960--1962 & \emph{Rubricarum instructum} and the 1962 Missal & Introduced the four-class system and the immediate preconciliar calendar against which the later reform is often compared.\\
4 Dec. 1963 & Vatican II, \emph{Sacrosanctum Concilium} 102--111 & Required revision of the liturgical year, protection of Sunday and the Proper of Time, Christological ordering of saints' feasts, and a more selective universal sanctoral.\\
14 Feb. 1969 & Paul VI, \emph{Mysterii Paschalis}, and \emph{Calendarium Romanum} & Promulgated the Universal Norms and revised General Roman Calendar, effective 1 January 1970. The new system uses solemnities, feasts, memorials, and optional memorials rather than the 1960 classes.\\
1970; 1975 & First and second typical editions of the Roman Missal & Embodied and then consolidated the reformed calendar in the Missal's liturgical structure.\\
1971; 1987--1988 & U.S. Guadalupe inscriptions & Added Our Lady of Guadalupe first as an obligatory memorial, then raised the national celebration to a feast.\\
1 Jan. 1996; 8 Dec. 1998 & Immaculate Heart rank and collision notifications & Raised the Immaculate Heart from optional to obligatory memorial, then established that when it coincides with another obligatory memorial, both are optional memorials for that year.\\
22 Jan.; 5 July 1999 & All-Americas Guadalupe feast; U.S. Ascension norm & Established Guadalupe as a feast in the national calendars of the Americas and allowed U.S. ecclesiastical provinces, under the approved procedure, to transfer the Ascension to the Seventh Sunday of Easter.\\
20 Apr. 2000; 2002; 2008 & Third typical edition and emended printing & The third typical edition was issued as typical in 2000, published in 2002, and reprinted with emendations in 2008. It supplies the present Latin base; the emended printing is the immediate Latin control for the 2011 U.S. edition.\\
5 May 2000 & Decree \emph{Misericors et miserator} & Added “or of Divine Mercy” to the title of the Second Sunday of Easter while prescribing continued use of the existing Missal and Hours texts.\\
26 June 2002 & \emph{Ex uberi terra} & Inscribed St Pius of Pietrelcina as an obligatory memorial on 23 September.\\
28 Sept. 2002; 2008 & Guadalupe and St Juan Diego integration & Inscribed the optional memorials of St Juan Diego on 9 December and Our Lady of Guadalupe on 12 December; both were incorporated into the 2008 emended printing.\\
24 July 2010; 27 Nov. 2011 & Confirmed U.S. adaptations and proper; English third edition & Three confirmations covered the U.S. adaptations, proper-calendar changes, and proper texts; the English edition came into use on 27 November 2011. Later national additions remain separate amendments.\\
12 Oct. 2012; 10 July 2013; 25 July 2014 & U.S. additions after the third edition & Added the U.S. optional memorials of St John Paul II, St Marianne Cope, and Blessed Francis Xavier Seelos. John Paul II's 2014 universal inscription later superseded the distinct national insertion.\\
29 May 2014 & Decree \emph{Pastor aeternus} & Added optional memorials of St John XXIII (11 October) and St John Paul II (22 October).\\
3 June 2016 & Decree on St Mary Magdalene & Raised 22 July from obligatory memorial to feast.\\
11 Feb.; 24 Mar. 2018 & Decree and notification on Mary, Mother of the Church & Added an obligatory memorial on the Monday after Pentecost and directed that it prevails when coinciding with another memorial of a saint or blessed.\\
25 Jan. 2019 & Decree on St Paul VI & Added his optional memorial on 29 May.\\
30 Sept. 2019 & Francis, \emph{Aperuit illis} & Designated the existing Third Sunday in Ordinary Time as the Sunday of the Word of God; the act prescribes no replacement rank or lectionary.\\
7 Oct. 2019 & Decree on Our Lady of Loreto & Added the optional memorial on 10 December.\\
18 May 2020 & Decree on St Faustina Kowalska & Added her optional memorial on 5 October.\\
25 Jan. 2021 & Three Doctors decree & Added optional memorials of St Gregory of Narek (27 February), St John of Avila (10 May), and St Hildegard of Bingen (17 September).\\
26 Jan. 2021 & Martha, Mary, and Lazarus decree & Replaced the memorial of St Martha alone on 29 July with the obligatory memorial of Saints Martha, Mary, and Lazarus.\\
21 Jan. 2022 & St Irenaeus declared Doctor of the Church & Added the title Doctor of the Church to the universal memorial's identifier without changing its 28 June date or memorial rank.\\
24 Dec. 2024; published 11 Feb. 2025 & Decree on St Teresa of Calcutta & Added her optional memorial on 5 September.\\
9 Nov. 2025; published 3 Feb. 2026 & Decree on St John Henry Newman & Added his optional memorial on 9 October; official 2026 calendars implement the entry.\\
17 July 2026 & Cutoff of this inventory & No later universal or United States amendment is asserted here.\\
\bottomrule
\end{longtable}
